% =========================================================================
% English translation (generated by AI using Claude Opus 4.8 (Max effort)).
% This file was translated from Hungarian to English by AI.
% DISCLAIMER: Please review against the original Hungarian .tex files, before relying on it!
% SOURCE: 3.1 Lineáris egyenletrendszerek interációs módszerei.tex
% =========================================================================
\documentclass[margin=0px]{article}

\usepackage{listings}
\usepackage[utf8]{inputenc}
\usepackage{graphicx}
\usepackage{float}
\usepackage[a4paper, margin=0.7in]{geometry}
\usepackage{amsthm}
\usepackage{amssymb}
\usepackage{amsmath}
\usepackage{fancyhdr}
\usepackage{setspace}

\onehalfspacing

\newenvironment{tetel}[1]{\paragraph{#1 \\}}{}

\pagestyle{fancy}
\lhead{\it{CS BSc Final Exam Topics}}
\rhead{3.1. Iterative methods for systems of linear equations}

\title{\textbf{{\Large ELTE IK - Computer Science BSc} \vspace{0.2cm} \\ {\huge Final Exam Topics}} \vspace{0.3cm} \\ 3.1. Iterative methods for systems of linear equations}
\author{}
\date{}

\begin{document}
\maketitle

\begin{center}
\fbox{\parbox{0.92\textwidth}{\small \textit{Note: This document was translated from Hungarian to English by AI. Please verify the information against the original Hungarian source before relying on it.}}}
\end{center}

\subsection{Iterative methods for systems of linear equations}

We approximate the system of linear equations (SLE) with sequences of vectors, striving for as fast convergence as possible.
The essence of the iterative methods is the transformation $Ax = b \Longleftrightarrow x = Bx + c$. Such a form exists,
and indeed it is not unique but can take many forms, and the various transformations provide the various iterative methods.\\

\noindent \textbf{Definition (contraction)}: The function $F: \mathbb{R}^{n} \to \mathbb{R}^{n}$ is a contraction if
$\exists 0 \leq q < 1 : \forall x,y \in \mathbb{R}^{n}:$
\begin{displaymath}
    \Vert F(x) - F(y) \Vert \leq q \Vert x - y \Vert
\end{displaymath}.

\noindent We call the value $q$ the contraction coefficient.\\

\noindent \textbf{Banach fixed-point theorem}: Let $F: \mathbb{R}^{n} \to \mathbb{R}^{n}$ be a contraction with contraction
coefficient $q$. Then the following statements are true:

\begin{enumerate}
    \item	$\exists! x^{*} \in \mathbb{R}^{n}: x^{*} = f(x^{*})$. We say that $x^{*}$ is the
          fixed point of the function $f$.

    \item	for $\forall x^{(0)} \in \mathbb{R}^{n}$ initial value the sequence $x^{(k+1)} = f(x^{(k)})$ is convergent,
          and $\lim \limits_{k\to\infty} x^{(k)} = x^{*}$.

    \item	$\Vert x^{(k)} - x^{*}\Vert \leq \frac{q^{k}}{1-q} \Vert x^{(1)} - x^{0}\Vert$.

    \item	$\Vert x^{(k)} - x^{*}\Vert \leq q^{k} \Vert x^{(0)} - x^{*}\Vert$.
\end{enumerate}

Note that with the transformation $Ax = b \Longleftrightarrow x = Bx + c$ we have established the connection with the
Banach fixed-point theorem, since now we look for the fixed point of the function $F(x) = Bx + c$. In the above formulation
$B$ is called the iteration matrix.\\

\noindent \textbf{Theorem (sufficient condition for convergence)}: If for the iteration matrix $B$ of the SLE $\Vert B \Vert < 1$,
then the iteration $x^{(k+1)} := Bx^{(k)} + c$ started from any $x^{(0)}$ converges to the solution of the SLE $Ax = b$.\\

\noindent \textbf{Theorem (necessary and sufficient condition for convergence)}: The iteration
$x^{(k+1)} := Bx^{(k)} + c$ started from any $x^{(0)}$ converges to the solution of the SLE $Ax = b$ $\Longleftrightarrow$
$\varrho(B) < 1$, where $\varrho(B) = \max_{1 \leq i \leq n} |\lambda_{i}(B)|$ is the spectral radius of the matrix $B$.

\subsubsection{Jacobi iteration}

Consider the $L + D + U$ decomposition of the matrix $A \in \mathbb{R}^{n \times n}$, where $L$ is the strict lower part of the matrix, $U$ is the
strict upper part, and $D$ is the diagonal part. With its help let us construct the following transformation:

\begin{displaymath}
    Ax = b \Longleftrightarrow (L + D + U)x = b \Longleftrightarrow Dx = -(L+U)x+b \Longleftrightarrow
    x = -D^{-1}(L+U)x + D^{-1}b
\end{displaymath}

\noindent The iteration matrix of the Jacobi iteration is therefore $B_{J} = -D^{-1}(L + U)$, and the iteration itself is:

\begin{displaymath}
    x^{(k+1)} :=  -D^{-1}(L + U)x^{(k)} + D^{-1}b
\end{displaymath}

\noindent Written in coordinate form:

\begin{displaymath}
    x^{(k+1)}_{i} :=
    -\frac{1}{a_{ii}}
    \Bigg[
    \sum_{\substack{j=1\\ j \not = i}}^{n} a_{ij}x_{j}^{(k)} - b_{i}
    \Bigg]
    \; (i = 1, ..., n)
\end{displaymath}

\noindent \textbf{Theorem}: If the matrix A is strictly diagonally dominant by rows, then
$\Vert B_{J} \Vert_{\infty} < 1$ (that is, the method is convergent).\\

\noindent \textbf{Theorem}: If the matrix A is strictly diagonally dominant by columns, then
$\Vert B_{J} \Vert_{1} < 1$ (that is, the method is convergent).

\subsubsection{Damped Jacobi iteration}

We still deal with the Jacobi iteration, but try to refine the method by introducing an extra parameter $\omega$.
Consider the equation $Dx = -(L + U)x + b$, as well as the trivial equation $Dx =
    Dx$. Multiply these respectively by $\omega$ and $1 - \omega$, then add the two
equations:

\begin{displaymath}
    Dx = (1 - \omega)Dx -\omega(L+U)x + \omega b
\end{displaymath}

\noindent Multiply by $D^{-1}$:

\begin{displaymath}
    x = (1 - \omega)Ix -\omega D^{-1}(L+U)x + \omega D^{-1}b \Longleftrightarrow
    x = ((1 - \omega)I -\omega D^{-1}(L+U))x + \omega D^{-1}b
\end{displaymath}

\noindent Based on this, $B_{J(\omega)} = (1 - \omega)I -\omega D^{-1}(L+U)$ and $c_{J}(\omega) = \omega D^{-1}b$.\\

\noindent Note that for $\omega = 1$ we get back exactly the Jacobi iteration.\\

\noindent Written in coordinate form:

\begin{displaymath}
    x^{(k+1)}_{i} =
    (1 - \omega)x^{(k)}_{i}
    -\frac{\omega}{a_{ii}}
    \Bigg[
    \sum_{\substack{j=1\\ j \not = i}}^{n} a_{ij}x_{j}^{(k)} - b_{i}
    \Bigg]
    \; (i = 1, ..., n)
\end{displaymath}

\noindent \textbf{Theorem}: If $J(1)$ is convergent, then for $\omega \in (0,1)$ so is $J(\omega)$.

\subsubsection{Gauss--Seidel iteration}

\noindent The idea of constructing another possible iteration is the following:

\begin{displaymath}
    Ax = b \Longleftrightarrow (L + D + U)x = b \Longleftrightarrow (L+D)x = -Ux+b \Longleftrightarrow
    x = -(L+D)^{-1}Ux + (L+D)^{-1}b
\end{displaymath}

\noindent This idea gives birth to the Gauss--Seidel iteration, that is:

\begin{displaymath}
    x^{(k+1)} :=  -(L+D)^{-1}Ux^{(k)} + (L+D)^{-1}b
\end{displaymath}

The iteration matrix of the iteration is therefore $B_{S} = -(L+D)^{-1}U$. To write the coordinate form we rewrite
the iteration a little:

\begin{displaymath}
    (L+D)x^{(k+1)} = -Ux^{(k)} + b
\end{displaymath}

\begin{displaymath}
    Dx^{(k+1)} = -Lx^{(k+1)} - Ux^{(k)} + b
\end{displaymath}

\begin{displaymath}
    x^{(k+1)} = -D^{-1} \big[Lx^{(k+1)} + Ux^{(k)} - b \big]
\end{displaymath}

\begin{displaymath}
    x^{(k+1)}_{i} =
    -\frac{1}{a_{ii}}
    \Bigg[
    \sum_{j=1}^{i-1} a_{ij}x_{j}^{(k+1)} +
    \sum_{j=i+1}^{n} a_{ij}x_{j}^{(k)}	-
    b_{i}
    \Bigg]
    \; (i = 1, ..., n)
\end{displaymath}

\noindent \textbf{Remark}: During implementation it is enough to store a single vector $x$ and overwrite its components in order, since
we can see that the first $i-1$ components are already taken from the ``new'' vector $x^{(k+1)}$.\\

\noindent \textbf{Theorem}: If $A$ is strictly diagonally dominant
\begin{enumerate}
    \item	by rows, then
          $\Vert B_{S} \Vert_{\infty} \leq \Vert B_{J} \Vert_{\infty} < 1$.

    \item  by columns, then
          $\Vert B_{S} \Vert_{1} \leq \Vert B_{J} \Vert_{1} < 1$.
\end{enumerate}

\noindent That is, Gauss--Seidel is also convergent, and at least as fast as Jacobi.

\subsubsection{Relaxation method}

The relaxation method essentially means the damped Gauss--Seidel iteration. To construct it,
consider the equations $(L+D)x = -Ux + b$ and $Dx = Dx$. Multiply these respectively by $\omega$ and
$1 - \omega$, then add the two equations:

\begin{displaymath}
    (D+\omega L)x = (1 - \omega)Dx -\omega Ux + \omega b
\end{displaymath}

\begin{displaymath}
    x = (D+\omega L)^{-1}\big[(1 - \omega)D -\omega U \big]x + \omega (D+\omega L)^{-1}b
\end{displaymath}

The iteration is therefore: $x^{(k+1)} = (D+\omega L)^{-1}\big[(1 - \omega)D -\omega U \big]x^{(k)} + \omega (D+\omega L)^{-1}b$, where
the iteration matrix is: $B_{S(\omega)}= (D+\omega L)^{-1}\big[(1 - \omega)D -\omega U \big]$. To write the coordinate form
we again rewrite the iteration a little:

\begin{displaymath}
    (D + \omega L)x^{(k+1)} = (1 - \omega)Dx^{(k)} -\omega Ux^{(k)} + \omega b
\end{displaymath}

\begin{displaymath}
    Dx^{(k+1)} =  -\omega Lx^{(k+1)} -\omega Ux^{(k)} + \omega b + (1 - \omega)Dx^{(k)}
\end{displaymath}

\begin{displaymath}
    x^{(k+1)} = -\omega D^{-1} \big[Lx^{(k+1)} + Ux^{(k)} - b \big] + (1 - \omega)x^{(k)}
\end{displaymath}

\begin{displaymath}
    x^{(k+1)}_{i} =
    -\frac{\omega}{a_{ii}}
    \Bigg[
    \sum_{j=1}^{i-1} a_{ij}x_{j}^{(k+1)} +
    \sum_{j=i+1}^{n} a_{ij}x_{j}^{(k)}	-
    b_{i}
    \Bigg] +
    (1 - \omega) x_{i}^{(k)}
    \; (i = 1, ..., n)
\end{displaymath}

\noindent Note that for $\omega = 1$ we get the Gauss--Seidel iteration.\\

\noindent \textbf{Theorem}: If the relaxation method is convergent started from every initial vector, then $\omega \in (0,2)$.\\

\noindent \textbf{Remark}: If $\omega \notin (0,2)$, then the method is generally not convergent (although for a given problem it may happen
that we find an initial vector from which the method converges).\\

\noindent \textbf{Theorem}: If $A$ is symmetric and positive definite and $\omega \in (0, 2)$, then the relaxation method is convergent. A
consequence of this is the convergence of the Gauss--Seidel iteration (the $\omega = 1$ case).\\

\noindent \textbf{Theorem}: If $A$ is tridiagonal, then $\varrho(B_{S}) = \varrho(B_{J})^{2}$, that is, the Jacobi and Gauss--Seidel iterations are
convergent or divergent at the same time.\\

\noindent \textbf{Theorem}: If $A$ is symmetric, positive definite and tridiagonal, then $J(1)$, $S(1)$ and $S(\omega)$ are convergent for $\omega \in (0, 2)$,
and for $S(\omega)$ the value of the optimal parameter is:

\begin{displaymath}
    \omega_{0} = \frac{2}{1 + \sqrt{1 - \varrho(B_{J})^{2}}}
\end{displaymath}

\subsubsection{Richardson iteration}

Let $p \in \mathbb{R}$. Thus

\begin{displaymath}
    Ax = b \Longleftrightarrow
    0 = -Ax + b \Longleftrightarrow
    0 = -pAx + pb \Longleftrightarrow
    x = (I-pA)x +pb
\end{displaymath}

The iteration is therefore $x^{(k+1)} := (I-pA)x^{(k)} +pb$. The iteration matrix: $B_{R(p)} = I-pA)$. The
vector $r^{(k)} := b - Ax^{(k)} $ is called the residual vector, since

\begin{displaymath}
    x^{(k+1)} = x^{(k)} - pAx^{(k)} + pb = x^{(k)} + pr^{(k)}
\end{displaymath}

\begin{displaymath}
    r^{(k+1)} = b - Ax^{(k+1)} = b - A(x^{(k)} + pr^{(k)}) = r^{(k)} - pAr^{(k)}
\end{displaymath}

\noindent Consider the algorithm of the construction: $r^{(0)} := b - Ax^{(0)}$, and because of the above:

\begin{displaymath}
    x^{(k+1)} := x^{(k)} + pr^{(k)}
\end{displaymath}

\begin{displaymath}
    r^{(k+1)} := r^{(k)} - pAr^{(k)}
\end{displaymath}

\noindent \textbf{Theorem}: If $A$ is symmetric, positive definite, and its eigenvalues are the following:
\begin{displaymath}
    0 < m := \lambda_{1} \leq \lambda_{2} \leq ... \leq \lambda_{n} =: M
\end{displaymath}

\noindent then for $p \in (0,\frac{2}{M})$ $R(p)$ is convergent, and the optimal parameter is: $p_{0} = \frac{2}{m + M}$.
Furthermore it is true that: $\varrho(B_{R(p_{0})}) = \frac{M - m}{M + m}$.

\subsection{Spline interpolation}

In the interpolation methods mentioned so far we worked with polynomials. It is also possible to try to fit another
type of function to the given point system. Piecewise polynomial functions that also satisfy certain continuity prescriptions,
the splines, have very advantageous properties.\\

\noindent \textbf{Spline of degree $l$}: Let a partition $\Omega_{n} = \left\{x_{0}, x_{1}, ..., x_{n}\right\}$ of the interval $[a,b]$ be given,
where $x_{0}=a, x_{n}=b$ and $l \in \mathbb{N}$. The function $s:[a,b] \to \mathbb{R}$ is a spline of degree $l$ with respect to
$\Omega_{n}$ if:

\begin{enumerate}
    \item	$s_{|[x_{k-1},x_{k}]}$ is a polynomial of degree $l$ for $\forall k = 1,..,n$

    \item	$s \in C^{l-1}[a,b]$, that is, on the whole interval it is $(l-1)$-times continuously differentiable
\end{enumerate}

\noindent Notation: $s_{l}(\Omega_{n})$ is the set of splines of degree $l$ belonging to $\Omega_{n}$.\\

\noindent \textbf{Spline interpolation}: Let the values $x_{k},f(x_{k})$ for $k=0,1,..,n$ and $l \in \mathbb{N}$ be given.
We look for the spline $s \in s_{l}(\Omega_{n})$ for which $s(x_{k}) = f(x_{k})$. For this we have to construct a
polynomial of degree $l$ for each interval. If we represent the polynomials by their coefficients, then this is $n(l+1)$ unknowns.
The number of prescribed conditions: $2n$ interpolation and $(l-1)(n-1)$ continuity conditions, since we only need to prescribe
the appropriate continuity condition for the respective derivatives at the interior points. The total number of conditions thus obtained
is $(l+1)n - (l-1)$, so for uniqueness $l-1$ conditions are still missing. These are given by so-called boundary conditions.
E.g. for cubic spline interpolation 2 boundary conditions are needed. These are the following (it is enough to choose one of them,
since each contains 2 conditions):

\begin{enumerate}
    \item	Natural boundary condition: $s''(a) = s''(b) = 0$.

    \item	Hermite boundary condition: $s'(a) = s_{a}, s'(b) = s_{b}$, where $s_{a}, s_{b}$ are prescribed numbers.

    \item	Periodic boundary condition (here we assume that $s(a) = s(b)$ also holds):
          $s'(a) = s'(b)$ and $s''(a) = s''(b)$
\end{enumerate}

\noindent \textbf{Construction of the first-degree spline}: The construction of the first-degree spline is trivial, by piecewise
linear Lagrange interpolation.\\

\noindent \textbf{Construction of the second-degree spline}: A single boundary condition is needed; let it be the following: $s'(a) = s_{a}$
for some number $s_{a}$. On the segment $[x_{0},x_{1}]$ we construct by Hermite interpolation the second-degree polynomial $H_{2}$
that satisfies the interpolation conditions and the boundary condition. The derivative of the resulting polynomial at $x_{1}$
is determined, so because of the continuous differentiability the value of the derivative at the left endpoint is given on the segment
$[x_{1},x_{2}]$. Applying Hermite interpolation again, we obtain the polynomial belonging to the segment $[x_{1},x_{2}]$.
Repeating this procedure we can construct the second-degree interpolation spline.\\

\noindent \textbf{Support of a function}: The set $supp(f) := \overline{\left\{x \in \mathbb{R}: f(x) \not = 0 \right\}}$ is called
the support of the function $f$.\\

\noindent \textbf{Partition of the number line}: $\Omega_{\infty} := \left\{...,x_{-2},x_{-1},x_{0},x_{1},...x_{n},x_{n+1},...\right\}$\\

\noindent \textbf{B-spline}: The system of spline functions $B_{l,k}, k \in \mathbb{Z}$ of degree $l$ is called the system of B-spline functions
if the following conditions hold:
\begin{enumerate}
    \item	$supp(B_{l,k}) = [x_{k},x_{k+l+1}]$, that is, its support is minimal

    \item	$B_{l,k}(x) \geq 0$

    \item	$\displaystyle\sum_{k \in \mathbb{Z}}B_{l,k}(x) = 1$
\end{enumerate}

\end{document}
